\documentclass{article}

\usepackage[english]{babel}
\usepackage{amsmath}
\usepackage{amssymb}
\usepackage{amsthm}
\usepackage[letterpaper,top=2cm,bottom=2cm,left=3cm,right=3cm,marginparwidth=1.75cm]{geometry}
\usepackage{graphicx}
\usepackage[colorlinks=true, allcolors=blue]{hyperref}
\usepackage{fancyhdr}
\usepackage{tikz}
\usetikzlibrary{decorations.markings,calc}
\usepackage{tikz-cd}
\usepackage{quiver}
\usetikzlibrary{matrix}
\usepackage[most]{tcolorbox}
\usepackage{hyperref}
\usepackage{array}
\usepackage{colonequals}
\usepackage{todonotes}
\usepackage{theoremref}

\font\maljapanese=dmjhira at 2.5ex
\newcommand{\yo}{\textrm{\!\maljapanese\char"48}}

\newtheorem{theorem}{Theorem}[section]

\theoremstyle{definition}
\newtheorem{definition2}[theorem]{Definition}
\newtheorem{lemma}[theorem]{Lemma}
\newtheorem{corollary}[theorem]{Corollary}
\newtheorem{definition}[theorem]{Definition}
\newtheorem{example}[theorem]{Example}
\newtheorem{examples}[theorem]{Examples}


\theoremstyle{remark}
\newtheorem*{remark}{Remark}

\theoremstyle{plain}
\newtheorem{proposition}[theorem]{Proposition}
\newtheorem{conjecture}[theorem]{Conjecture}


\newcommand{\R}{\mathbb{R}}
\newcommand{\C}{\mathbb{C}}
\newcommand{\Z}{\mathbb{Z}}
\newcommand{\N}{\mathbb{N}}
\newcommand{\Q}{\mathbb{Q}}
\newcommand{\mb}[1]{\mathbb{#1}}
\newcommand{\mc}[1]{\mathcal{#1}}
\newcommand{\mk}[1]{\mathfrak{#1}}
\newcommand{\un}{\cup}
\newcommand{\ic}{\cap}
\pagestyle{fancy}
\newcommand\size{1}% distance of nodes from center

\usepackage{microtype}

\usepackage{caption}
\captionsetup[figure]{labelformat=empty}%

\begin{document}

\title{N\'eron Models of Elliptic Curves}

\author{Daniel Ao}
\maketitle

\tableofcontents

\section{Introduction}


\section{Group Schemes and Arithmetic Surfaces}

\subsection{Group Schemes}

We will see that the primary objects of interest when it comes to N\'eron models are group schemes.
In this section we introduce group schemes from the functor of points view, and also show that its equivalent to the traditional definitions.
At the end, we introduce the notion of arithmetic surfaces which we will later use to show the existence of N\'eron models for elliptic curves over Dedekind domains.

\begin{definition}
	An $S$-scheme $G$ is a group scheme if the functor 
	\[\text{Hom}(-, G): \textbf{Sch/S} \to \textbf{Set}\]
 factors through the category $\textbf{Group}$.
\end{definition}

In other words, for any $S$-scheme $T$, the set of $T$-points of $G$ forms a group, which is functorial in $T$.
We show that this is equivalent to the usual definitions of groups schemes, being the existence of maps 
\[e: \text{Spec}(S) \to G, \quad \quad i:G \to G, \quad \quad \mu: G \times G \to G, \]
such that the following squares commute:

\todo{Proof of equivalence}

\begin{example}
	An elliptic curve is a group scheme.
\end{example}

\begin{definition}
	Abelian schemes/varieties.
\end{definition}

\begin{remark}
	There are some benefits to the functor of points view. 
	For example, from our definition it becomes clear that a product of group schemes is still a group scheme. 
	As another aside, $\mathbb{A}^1$ represents the global section functor has a natural ring structure, the same arguments
\end{remark}

The multiplication by $m$ endomorphisms are very important in the study of elliptic curves. 
In fact, these are defined for any abelian group scheme.

\begin{definition}
The multiplication by $m$ endomorphism map for a group scheme $X$ over $S$, is the map induced by Yoneda from the natural transformation $(-)^n \,:\, \text{Hom}(-,X) \to \text{Hom}(-,X)$, which for each $S$-scheme $U$, the component $\text{Hom}(U,X) \to \text{Hom}(U,X)$ is the multlpication by $n$ mp.
\end{definition}

Similarly, we can define the "translation by $x$" map, for $x$ an $X$-point of

\begin{definition}
	
\end{definition}


\subsection{Arithmetic Surfaces}


\begin{definition}
	Let $C/K$ be a non-singular connected projective curve. 
	An \textit{arithmetic surface} (over $R$) is an integral, normal, excellent scheme that is flat and finite type over $\text{Spec}(R)$ whose generic fiber is isomorphic to $C/K$ and whose special fibers are unions of curves.
\end{definition}


\begin{example}
$\mb{P}^1_R$  perhaps?
\end{example}

\begin{proposition}
	Let $\pi: C \to \text{Spec}(R)$ be a regular arithmetic surface, and let $\mk{p} \in \text{Spec}(R)$.
	Then
\begin{enumerate}
	\item Let $x \in C_{\mk{p}} \subset C$ a closed point on a fiber of $C$ over $\mk{p}$.
		Denote $\pi^*: R_{\mk{p}} \to \mc{O}_{C,x}$ the map of stalks induced by $\pi$ and $\mk{m}_{C,x}$ the maximal ideal of $\mc{O}_{C,s}$.
		Then 
		\[\text{$C_{\mk{p}}$ is non-singular at $x$} \Longleftrightarrow \pi^*(\mk{p}) \subsetneq \mk{m}^2_{C,x}\]
	\item Let $P \in C(R)$, then $C_{\mk{p}}$ is non-singular at $P(\mk{p})$. 
\end{enumerate}

\end{proposition}


\begin{proof}
	\cite{AEC2}, pg.327
\end{proof}

\begin{proposition}
	Let $X/\text{Spec}(R)$ be an arithmetic surface, and $C/K$ its generic fiber.
\begin{enumerate}
	\item if $C$ is proper over $\text{Spec}(R)$, then 
		\[C(K) = X(R)\]
	\item Suppose furthermore that $X$ is regular, and $X^0 \subset X$ the largest open subscheme of $X$ that is smooth over $\text{Spec}(R)$. 
		Then 
		\[X(R) = X^0(R)\]
\end{enumerate}

\begin{proof}
	\cite{AEC2}, pg.327
\end{proof}

\end{proposition}

\begin{theorem}[Minimal Proper Regular Models (Abhyankar, Lipman, Lichtenbaum, Shafarevich)]
	Let $C/K$ be a non-singular projective curve of genus $g$.
	Then there exists a proper, regular arithmetic over $R$ with generic fiber isomorphic to $C/K$ which we call proper regular models for $C/K$.
	Furthermore, if $g \geq 1$, then there exists a minimal proper regular model $C^{min}/K$ satisfying the following minimality property: \newline
	\indent \textbf{Let $C'/R$ be any other proper regular model for $C/K$. Then any $R$-birational map $C' \to C^{min}$ inducing an isomorphism on the generic fiber is an $R$-isomorphism.} 
\end{theorem}


\section{N\'eron Models}

Throughout this section, let $S = \text{Spec}(R)$ denote the spectrum of a Dedekind domain $R$ with fraction field $K$, for example let $R$ be the ring of integers of a number field. 
For any scheme $Y$ over $S$, we will denote $Y_K$ for its generic fiber $Y \times_{S} \text{Spec}(K)$.
Most of this section can be done more generally for Dedekind Schemes (noetherian normal schemes of dimension $\leq 1$), where we take $K = \prod_i K(\eta_i)$ for $\eta_i$ ranging over the generic components of the (finitely many) irreducible components.
As such many theorems will use this terminology. \newline

\indent As 
\todo{Dedekind Schemes?}



Given a scheme $X$ over Spec($K$), our interest will be in "lifts" of $X$ to a scheme $X'$ over $S$ such that its generic fiber $X'_K$ is isomorphic to $X$, which we call $S$-models of $X$.  
One naturally asks if there are minimal such lifts, which gives us the most general definition of a Neron Model, as in \cite{BLR}.

\begin{definition}
	Let $X_K$ be a smooth and separated $K$-scheme of finite type. A N\'eron model of $X_K$ is an $S$-model of $X$  of $X_K$ which is smooth, separated, and of finite type which satisfies the following universal property (which we call the N\'eron mapping property).\\
	\indent \textbf{For each smooth $S$-scheme $Y$ and $K$-morphism $u_K: Y_K \to X_K$, there exists a unique $S$-morphism $\mu: Y \to X$ extending $u_K$}.
\end{definition}

\noindent Immediately from its definition we have the following important facts.

\begin{proposition}
Let $X$ be a N\'eron model for its generic fiber $X_K$. Then the following hold:
\begin{enumerate}
  \item $X$ is uniquely determined by $X_K$ up to unique isomorphism.
  \item The natural map $X(S) \to X_K(K)$ is bijective, in other words every $K$-point of $X_K$ lifts to an $S$-point of $X$. 
\end{enumerate}
\end{proposition}

\begin{proof}
	The proof of (a) is the usual one for objects defined by universal properties. Namely, for two such models $X$ and $X'$ the identity morphisms  on their generic fibers induce maps $f:X \to X'$ and $g: X' \to X$ from the N’eron mapping property. Then $f \circ g$ and $g \circ f$ are unique lifts of the identity morphisms of the generic fibers, but the identity morphism already satisfies this so they must be equal. 
	Hence $f \circ g = \text{id}_{X'}$ and $g \circ f = \text{id}_{X}$ which shows $X$ and $X'$ are isomorphic. \newline
\indent
\todo{Complete proofs}
\end{proof}

\subsection{N\'eron Models of Group Schemes}

\todo{intro..}


\begin{proposition}
	Let $X$ be a smooth and separated $S$-scheme which is a N\'eron model of its generic fiber $X_K$.
	If $X_K$ is a $K$-group scheme, then the group scheme structure extends to an $S$-group scheme structure on $X$.
\end{proposition}

\begin{proposition}
	Let $S$ be a normal noetherian base scheme, and $u: Z \to G$ an $S$-rational map from a smooth $S$-scheme $Z$ and separated $S$-group shceme $G$.
	If $u$ is defined in codimension $\leq 1$, it is defined everywhere.
\end{proposition}

\begin{proof}
	\cite{BLR}, pg.116, fill in and do, might need a lemma.
\end{proof}

\begin{proposition}
	Let $X$ be a proper and smooth $S$-group scheme with connected fibers. 
	Then $X$ is a N\'eron model of its generic fiber $X_K$.
\end{proposition}

\begin{proof}
	\cite{BLR}, pg.22
\end{proof}

\subsection{Basic Properties of N\'eron Models}

\begin{proposition}
	N\'eron models commute with \'etale base change, i.e. if $S' \to S’$ is  \'etale, with $K'$ the ring of rational functions on $S'$, then $X \times_S S'$ is a N\'eron model for $X_K \times_{\text{Spec(K)}} \text{Spec}(K')$. 
\end{proposition}

\begin{proof}
	Basic diagram chasing/properties, TODO.	
\end{proof}

As an example, as open immersions are \'etale, for any open subscheme $S' \subset S$ we have that $X \otimes_S S'$ is a N\'eron model over $S'$.
So if 

\begin{proposition}
	Let $s$ be a point of $S$. Then
\begin{enumerate}
	\item If $X,Y$ are $S$-schemes of finite presentation, then the canonical map
		\[\varinjlim \text{Hom}_{S'}(X \times_S S'. Y \times_S S') \to \text{Hom}_{\mc{O}_{S,s}}(X \otimes_S \mc{O}_{S,s}, Y \otimes_S \mc{O}_{S,s}) \]
	taken over all open neighborhoods $S'$ of $s$ in $S$ is bijective.
\item Let $X'$ be an $\mc{O}_{S,s}$-scheme of finite presentation.
	Then there is an open neigborhood $S'$ of $s$ in $S$ and an $S'$-scheme $X'$ of finite presentation such that $X \otimes_{S'} \mc{O}_{S,s}$ is isomorphic to $X'$.
\end{enumerate}

\end{proposition}

\begin{proposition}
	Let $S$ be a Dedekind Scheme and $X$ an $S$-scheme of finite type.
	Then the following are equivalent
	\begin{enumerate}
		\item $X$ is a N\'eron model of its generic fiber
		\item For each closed point $s \in S$, the $\mc{O}_{S,s}$-scheme $X \times_S \mc{O}_{S,s}$ is a N\'eron model of its generic fiber.
	\end{enumerate}
\end{proposition}

\begin{proof}
	\cite{BLR}, pg.21, uses an EGA (prev) lemma.
\end{proof}

\begin{proposition}
	Let $S$ be a Dedekind scheme with ring of rational functions $K$, and let $X_K$ be a smooth and separated $K$-scheme of finite type.
	Then the following are equivalent.
\begin{enumerate}
	\item There exists a global N\'eron model $X$ of $X_K$ over $S$.
	\item There exists a dense open subscheme $S' \subset S$ such that $X_K$ admits a N\'eron model over $S'$ as well as N\'eron models for the (finitely many) closed points of $S - S'$.
\end{enumerate}
\end{proposition}

\begin{proof}
	\cite{BLR}, pg.25, also uses EGA Lemma.
\end{proof}



\section{Elliptic Curves}

In this section, we prove the existence of N\'eron models for elliptic curves over $K$, the fraction field of a Dedekind domain $R$. 
The majority of our time will be spent in showing the existence for the local case where $R$ is a discrete valuation rings, which will requires the introduction of the notion of special henselization. \newline

Before we get too deep, let's just consider the problem of constructing $S$-models of $K$-schemes without worrying about them having nice properties quite yet.
For simplicity, assume that $S = \text{Spec}(R)$ for $R$ a Dedekind domain.
In the affine case, suppose we have $X \cong \text{Spec}(K[x_1, \dots, x_n]/I)$, where we may let $I = (f_1, \dots, f_n)$ have finitely many generators by noetherianity. 
We do the most natural thing and note that we can scale the finitely many $f_i$ by a nonzero common multiple (though not unique) $r \in R$ such that each $rf_i \in R[x]$.
In this way, we have that $R[x]/(rf_1, \dots, rf_n)$ is a model for $X$ as the generic fiber is the spectrum of 
\[K \otimes_K R[x]/(rf_1, \dots, rf_n) = K[x]/(rf_1, \dots, rf_n) = K[x]/(f_1, \dots, f_n)\]
as $r$ is a unit in $K$. \newline

\indent Similarly, in the projective case when $K \cong \text{Proj}(K[x_1, \dots, x_n]/(g_1, \dots, g_n))$ for homogeneous $g_i$ we can similarly scale them by nonzero (and still non-unique) $r \in R$ to obtain integral models $\text{Proj}(R[x_1, \dots, x_n]/(rg_1, \dots, rg_n))$  for $X$.
In fact, note that $X$ is still projective, so satisfies the property that $K$-points extend uniquely to $R$-points, an important consequence of the N\'eron mapping property ().



An elliptic curve $E/K$ is a projective curve, so as mentioned in the discussion above there are projective $\text{Spec}(R)$-models for $E$ which are potential candidates for N\'eron models.
In this section we will show that in the case when $E$ has good reduction the N\'eron mdoel can be recovered this way.
However, more work will need to be done when there is bad reduction.

\indent First 


\begin{theorem}
	Let $E/K$ an elliptic curve.
	Choose a Weierstrass equation 
	\[y^2z + a_1xyz + a_3yz^2 = x^3 + a_2x^2z + a_4xz^2 + a_6z^3\]
	with coefficients in $R$ which defines a projective subscheme of $W \subset \mb{P}^2_R$.
	Let $W^0$ be the largest open subscheme of $W$ which is smooth over $R$. 
	Then 
\begin{enumerate}
	\item Both $W/R$ and $W^0/R$ have generic fiber $E/K$.
	\item The natural map $W(R) \to E(K)$ is a bijetion. 
		Furthermore, if $R$ is regular then the natural map $W^0(R) \to W(R)$ is also also a  bijection.
	\indent The addition and negation maps on $E$ extend to $R-$morphisms
	\[W^0(R) \times W^0(R) \to W^0(R), \quad \quad W^0(R) \to W^0(R) \]
	which make $W^0$ into a group scheme over $R$.
\end{enumerate}

\end{theorem}

\begin{proof}
	Uses 4.4(a), (b) from \cite{AEC2}, pg.329. 
\end{proof}

\begin{theorem}
	Let $E/K$ an elliptic curve.
	Choose a Weierstrass equation 
	\[y^2z + a_1xyz + a_3yz^2 = x^3 + a_2x^2z + a_4xz^2 + a_6z^3\]
	with coefficients in $R$ which defines a projective subscheme of $W \subset \mb{P}^2_R$.
	If $W$ is smooth over $R$, or equivalently the Weierstrass equation has good reduction at every prime of $R$, then $W/R$ is a N\'eron model for $E/K$.
\end{theorem}

\begin{proof}
	Previous theorem says it extends to a groups cheme, now use Proposition 3.5..
\end{proof}

\subsection{The Local Case}

\begin{definition}[(Strictly) Henselian Local Rings]
	A discrete valuation ring $R$ is called \textit{Henselian} if it satisfies Hensel's lemma.
	That is, for any monic polynomial $f(x) \in R[x]$ and $a \in R$ satisfying
	\[f(a) \equiv a \pmod{\mk{p}} \quad \quad f'(a) \not\equiv 0 \pmod{\mk{p}}\]
	there exists a unique element lift of $a$ to  a $\alpha \in R$ such that $f(\alpha) = 0$.
	Furthermore, if $R$ has separably closed residue field it is called \textit{strictly Henselian}.
\end{definition}

\begin{proposition}
	Let $(R, \mk{m}, k)$ be a discrete valuation ring. 
	Let $X/R$ be a smooth $R$-scheme and $\overline{X}/R$ it special fiber. 
	If $R$ is Henselian, then the reduction map
	\[X(R) \to \overline{X}(k)\]
	is surjective.
\end{proposition}

\begin{proof}
	\cite{AEC2}, pg 342.
\end{proof}


\begin{proposition}[(Strict) Henselization of DVRs]
	Let $(R, \mk{m}, k)$ be a discrete valuation ring with fraction field $K$.
	Denoting $K^s$ the separable closure and $R^s$ the integral closure of $R$ in $K^s$. 
	Let $\mk{p}$ be an ideal lying above $\mk{m}$ and $D$ and $I$ the corresponding decomposition and inertia groups.
	\begin{enumerate}
		\item Let $R^s(D)$ denote the subring of $R^s$ fixed by $D$ and define $R^h$ the lcoalization of $R^s(D)$ at the maximal ideal $\mk{p} \ic R^s(D)$.
		Then $R^h$ is Henselian and is called the \textit{Henselization} of $R$.
	\item Let $R^s(I)$ denote the subring of $R^{s}$ fixed by $D$ and define $R^{sh}$ the lcoalization of $R^s(I)$ at the maximal ideal $\mk{p} \ic R^s(I)$.
		Then $R^{sh}$ is Henselian and is called the \textit{strict Henselization} of $R$.
	\end{enumerate}
	
\end{proposition}

\subsection{The Global Case}

\begin{theorem}
	Let $E/K$ be an elliptic curve, and $X/R$ a minimal proper regular model for $E/K$, and $X^0/R$ the largest open subscheme of $X/R$ which is smooth over $R$.
	Then $X^0/R$ is a N\'eron model for $E/K$/
\end{theorem}

\begin{proof}
	
\end{proof}


\section{Sample Computations of N\'eron Models}



\bibliographystyle{plain}
\bibliography{lib}

\end{document}
